\section{XÁC SUẤT CỔ ĐIỂN}
\subsection{\vdmh}
\Opensolutionfile{ans}[ans/DA-CD5-XAC-SUAT-CO-DIEN-VD]
\begin{vd}%[135941] [MĐ1]
    Từ một hộp chứa $7$ quả cầu màu đỏ và $5$ quả cầu màu xanh, lấy ngẫu nhiên đồng thời $3$ quả cầu. Xác suất để lấy được $3$ quả cầu màu xanh bằng
    \choice
    {$\frac{5}{12}$}
    {$\frac{7}{44}$}
    {\True $\frac{1}{22}$}
    {$\frac{2}{7}$}
    \loigiai{
        Không gian mẫu: Lấy $3$ quả cầu từ $12$ quả cầu $\Rightarrow n(\Omega) = C_{12}^3 = 220$.\\
        Gọi biến cố $A$: "Lấy được $3$ quả cầu màu xanh".\\
        Số cách lấy $3$ quả màu xanh từ $5$ quả màu xanh là $n(A) = C_5^3 = 10$.\\
        Xác suất $P(A) = \frac{10}{220} = \frac{1}{22}$.
    }
\end{vd}
\begin{vd}%[792250] [MĐ2]
    Tổ một có $3$ học sinh nam và $4$ học sinh nữ. Tổ hai có $5$ học sinh nam và $2$ học sinh nữ. Chọn ngẫu nhiên mỗi tổ một học sinh đi làm nhiệm vụ. Tính xác suất sao cho chọn được hai học sinh có cả nam và nữ.
    \choice
    {$\frac{25}{49}$}
    {\True $\frac{26}{49}$}
    {$\frac{2}{13}$}
    {$\frac{5}{13}$}
    \loigiai{
        Không gian mẫu: $n(\Omega) = C_7^1 \cdot C_7^1 = 7 \cdot 7 = 49$.\\
        Các trường hợp để có cả nam và nữ:
        \begin{itemize}
            \item Nam ở tổ $1$ và nữ ở tổ $2$: $3 \cdot 2 = 6$ cách.
            \item Nữ ở tổ $1$ và nam ở tổ $2$: $4 \cdot 5 = 20$ cách.
        \end{itemize}
        $n(A) = 6 + 20 = 26$.\\
        Xác suất $P(A) = \frac{26}{49}$.
    }
\end{vd}
\begin{vd}
    Cho tập $S=\{1;2;3;4;5;6;7\}$. Chọn ngẫu nhiên một số tự nhiên có $5$ chữ số khác nhau được lập từ tập $S$. Tính xác suất để số được chọn có hai chữ số $1$ và $2$ đứng cạnh nhau? (làm tròn kết quả đến hàng phần trăm)\shortans{0,19}
    \loigiai{
        Số phần tử của không gian mẫu (số cách lập số có 5 chữ số khác nhau từ $S$) là: $n(\Omega) = A_7^5 = 2520$.\\
        Gọi biến cố $A$: "Số được chọn có hai chữ số $1$ và $2$ đứng cạnh nhau".\\
        Gọi số cần tìm có dạng $\overline{abcde}$. Gộp 1 và 2 thành một cụm $X$.
        \begin{itemize}
            \item Số cách hoán vị 1 và 2 trong $X$ là $2! = 2$ cách.
            \item Xếp $X$ và 3 chữ số khác lấy từ $S \setminus \{1,2\}$ (có 5 chữ số). Số cách chọn 3 chữ số từ 5 chữ số $\{3;4;5;6;7\}$ và xếp vào 3 vị trí (cùng với $X$ là 4 vị trí) là $A_5^3 \times 4 = 240$ cách.
        \end{itemize}
        Số các số thỏa mãn (kết quả thuận lợi cho $A$) là $n(A) = 2 \times 240 = 480$.\\
        Xác suất cần tìm là $P(A) = \frac{480}{2520} = \frac{4}{21} \approx 0,19$.
        \hfill \textbf{Đáp số:} $0,19$.
    }
    \dotlineEX{1}
\end{vd}

\begin{vd}
    Xếp ngẫu nhiên $5$ cuốn sách Toán khác nhau, $3$ cuốn sách Lí khác nhau và $4$ cuốn sách Hóa khác nhau lên một kệ dài. Biết xác suất để các cuốn sách Toán luôn đứng cạnh nhau và không có hai cuốn sách Lí nào đứng cạnh nhau là phân số tối giản $\frac{a}{b}$. Tính giá trị của $a+b$.\shortans{278}
%    \loigiai{
%        Tổng số sách là $5+3+4 = 12$ cuốn.\\
%        Không gian mẫu (xếp ngẫu nhiên 12 cuốn sách): $n(\Omega) = 12!$.\\
%        Gọi biến cố $A$: "Sách Toán đứng cạnh nhau và sách Lí không đứng cạnh nhau".
%        \begin{itemize}
%            \item Gộp 5 sách Toán thành 1 cụm $X$. Hoán vị 5 sách Toán trong $X$: $5!$ cách.
%            \item Xếp cụm $X$ và 4 sách Hóa: có $(1+4)! = 5!$ cách.
%            \item 5 phần tử trên tạo ra 6 khoảng trống. Xếp 3 sách Lí vào 6 khoảng trống để chúng không kề nhau: $A_6^3$ cách.
%        \end{itemize}
%        Số kết quả thuận lợi: $n(A) = 5! \times 5! \times A_6^3 = 120 \times 120 \times 120 = 1728000$.\\
%        Xác suất $P(A) = \frac{1728000}{12!} = \frac{1}{277}$.\\
%        Vậy $a=1, b=277 \Rightarrow a+b = 278$.
%        \hfill \textbf{Đáp số:} $278$.
%    }
\dotlineEX{4}
\end{vd}
\begin{vd}
    Tại một nút giao thông đông đúc, người ta quy định chu kỳ của đèn tín hiệu giao thông gồm $9$ pha liên tiếp, mỗi pha kéo dài $20$ giây. Trong mỗi chu kỳ phải có: $3$ pha đèn đỏ, $3$ pha đèn vàng và $3$ pha đèn xanh. Mỗi pha được sắp xếp ngẫu nhiên theo một thứ tự nhất định trong chu kỳ. Biết xác suất để trong một chu kỳ hai pha cùng màu không được xuất hiện liền nhau bằng $\frac{a}{b}$ với $a,b \in \mathbb{N}^*$ và $\frac{a}{b}$ là phân số tối giản. Khi đó $a+b$ bằng bao nhiêu?\shortans{309}
    \loigiai{
        Giả sử có 3 Đỏ (Đ), 3 Vàng (V), 3 Xanh (X).\\
        Tổng số cách xếp ngẫu nhiên 9 pha: $n(\Omega) = \frac{9!}{3!3!3!} = 1680$.\\
        Gọi biến cố $A$: "Không có 2 pha cùng màu xuất hiện liền nhau".\\
        Sử dụng nguyên lý bao hàm loại trừ (hoặc tra cứu/tính toán cụ thể bài toán xếp đồ vật có lặp), số cách xếp thỏa mãn là $n(A) = 174$.\\
        Xác suất $P(A) = \frac{174}{1680} = \frac{29}{280}$.\\
        Do đó $a=29, b=280 \Rightarrow a+b = 309$.
        \hfill \textbf{Đáp số:} $309$.
    }
\end{vd}
\begin{vd}
    Quanh một khu vực hình vuông, người ta lắp $20$ chiếc camera theo dõi cách đều nhau. Tuy nhiên vì mật độ quá dày nên người ta quyết định chọn ngẫu nhiên $5$ chiếc camera để tháo dỡ. Tính xác suất để $5$ chiếc camera được chọn tháo dỡ không có bất kỳ $2$ chiếc nào nằm cạnh nhau? (làm tròn kết quả đến hàng phần trăm)\shortans{0,26}
    \loigiai{
        Số phần tử của không gian mẫu (chọn ngẫu nhiên 5 camera từ 20 camera): $n(\Omega) = C_{20}^5 = 15504$.\\
        Gọi biến cố $A$: "5 chiếc camera được chọn không có 2 chiếc nào nằm cạnh nhau".\\
        Vì các camera cách đều nhau quanh hình vuông nên tính chất kề nhau tương đương với việc xếp trên một vòng tròn. Đây là bài toán chọn $k=5$ phần tử không kề nhau từ $n=20$ phần tử xếp trên vòng tròn.\\
        Số kết quả thuận lợi cho $A$ là (áp dụng công thức từ bài toán chia kẹo Euler): 
        $$n(A) = \frac{n}{n-k} C_{n-k}^k = \frac{20}{15} C_{15}^5 = 4004.$$
        Xác suất cần tìm: $P(A) = \frac{4004}{15504} = \frac{1001}{3876} \approx 0,26$.
        \hfill \textbf{Đáp số:} $0,26$.
    }
\end{vd}
\begin{vd}
    Có $12$ chiếc kẹo giống hệt nhau. Người ta chia ngẫu nhiên số kẹo này cho $3$ em bé A, B, C sao cho mỗi em nhận được ít nhất $1$ chiếc kẹo. Biết xác suất để số kẹo của em A nhận được là một số chia hết cho $3$ có dạng phân số tối giản $\frac{a}{b}$. Tính giá trị của biểu thức $a+b$.\shortans{14}
    \loigiai{
        Gọi số kẹo của 3 em bé A, B, C nhận được lần lượt là $x, y, z$ với $x,y,z \in \mathbb{N}^*$.\\
        Ta có phương trình: $x + y + z = 12$.\\
        Số phần tử của không gian mẫu (số nghiệm nguyên dương của phương trình): $n(\Omega) = C_{12-1}^{3-1} = C_{11}^2 = 55$.\\
        Gọi biến cố $E$: "Số kẹo của em A chia hết cho 3". Khi đó $x \vdots 3 \Rightarrow x \in \{3; 6; 9\}$.
        \begin{itemize}
            \item Nếu $x = 3 \Rightarrow y + z = 9 \Rightarrow$ Số nghiệm nguyên dương là $C_{9-1}^{2-1} = 8$.
            \item Nếu $x = 6 \Rightarrow y + z = 6 \Rightarrow$ Số nghiệm nguyên dương là $C_{6-1}^{2-1} = 5$.
            \item Nếu $x = 9 \Rightarrow y + z = 3 \Rightarrow$ Số nghiệm nguyên dương là $C_{3-1}^{2-1} = 2$.
        \end{itemize}
        Số kết quả thuận lợi cho biến cố $E$ là: $n(E) = 8 + 5 + 2 = 15$.\\
        Xác suất cần tìm là: $P(E) = \frac{15}{55} = \frac{3}{11}$.\\
        Suy ra $a=3, b=11 \Rightarrow a+b = 14$.
        \hfill \textbf{Đáp số:} $14$.
    }
\end{vd}
\Closesolutionfile{ans}
\subsection{\btrl}
\Opensolutionfile{ans}[ans/DA-CD5-XAC-SUAT-CO-DIEN-BTRL]
\begin{ex}%[45919] [MĐ1]
    Chọn ngẫu nhiên hai số khác nhau từ $23$ số nguyên dương đầu tiên. Xác suất để chọn được hai số có tổng là một số chẵn bằng
    \choice
    {\True $\frac{11}{23}$}
    {$\frac{1}{2}$}
    {$\frac{265}{529}$}
    {$\frac{12}{23}$}
    \loigiai{
        Tập hợp $23$ số nguyên dương đầu tiên có $12$ số lẻ và $11$ số chẵn.\\
        Không gian mẫu: $n(\Omega) = C_{23}^2 = 253$.\\
        Để tổng hai số là số chẵn thì hai số được chọn phải cùng chẵn hoặc cùng lẻ.\\
        Số cách chọn thỏa mãn: $n(A) = C_{12}^2 + C_{11}^2 = 66 + 55 = 121$.\\
        Xác suất $P(A) = \frac{121}{253} = \frac{11}{23}$.
    }
\end{ex}
\begin{ex}%[234260] [MĐ1]
    Chọn ngẫu nhiên một số từ tập hợp các số tự nhiên thuộc đoạn $[30; 50]$. Xác suất để chọn được số có chữ số hàng đơn vị lớn hơn chữ số hàng chục bằng
    \choice
    {\True $\frac{11}{21}$}
    {$\frac{8}{21}$}
    {$\frac{13}{21}$}
    {$\frac{10}{21}$}
    \loigiai{
        Tập hợp có $50 - 30 + 1 = 21$ số. Không gian mẫu $n(\Omega) = 21$.\\
        Các số có chữ số hàng đơn vị lớn hơn hàng chục là:
        \begin{itemize}
            \item Hàng chục là $3$: $34, 35, 36, 37, 38, 39$ (có $6$ số)
            \item Hàng chục là $4$: $45, 46, 47, 48, 49$ (có $5$ số)
            \item Hàng chục là $5$: không có.
        \end{itemize}
        Suy ra $n(A) = 6 + 5 = 11$.\\
        Xác suất $P(A) = \frac{11}{21}$.
    }
\end{ex}
\begin{ex}%[322507] [MĐ1]
    Một hộp chứa $11$ quả cầu gồm $5$ quả cầu màu xanh và $6$ quả cầu màu đỏ. Chọn ngẫu nhiên đồng thời $2$ quả cầu từ hộp đó. Xác suất để $2$ quả cầu chọn ra cùng màu bằng
    \choice
    {$\frac{5}{22}$}
    {$\frac{6}{11}$}
    {\True $\frac{5}{11}$}
    {$\frac{8}{11}$}
    \loigiai{
        Không gian mẫu: $n(\Omega) = C_{11}^2 = 55$.\\
        Hai quả cầu cùng màu có thể là cùng xanh hoặc cùng đỏ: $n(A) = C_5^2 + C_6^2 = 10 + 15 = 25$.\\
        Xác suất $P(A) = \frac{25}{55} = \frac{5}{11}$.
    }
\end{ex}
\begin{ex}%[792229] [MĐ1]
    Để kiểm tra chất lượng sản phẩm từ một công ty sữa, người ta đã gửi đến bộ phận kiểm nghiệm $5$ hộp sữa cam, $4$ hộp sữa dâu và $3$ hộp sữa nho. Bộ phận kiểm nghiệm chọn ra $3$ hộp sữa để phân tích mẫu. Tính xác suất để $3$ hộp sữa được chọn có đủ cả $3$ loại.
    \choice
    {\True $\frac{3}{11}$}
    {$\frac{6}{11}$}
    {$\frac{5}{11}$}
    {$\frac{8}{11}$}
%    \loigiai{
%        Tổng số hộp sữa là $12$. Không gian mẫu: $n(\Omega) = C_{12}^3 = 220$.\\
%        Để chọn $3$ hộp có đủ $3$ loại, ta chọn mỗi loại $1$ hộp: $n(A) = C_5^1 \cdot C_4^1 \cdot C_3^1 = 5 \cdot 4 \cdot 3 = 60$.\\
%        Xác suất $P(A) = \frac{60}{220} = \frac{3}{11}$.
%    }
\dotlineEX{1}
\end{ex}
\begin{ex}%[792230] [MĐ1]
    Từ một hộp chứa $16$ thẻ được đánh số từ $1$ đến $16$, chọn ngẫu nhiên ra $4$ thẻ. Tính xác suất để $4$ thẻ được chọn đều được đánh số chẵn.
    \choice
    {\True $\frac{1}{26}$}
    {$\frac{3}{26}$}
    {$\frac{1}{4}$}
    {$\frac{5}{9}$}
%    \loigiai{
%        Tập hợp $16$ thẻ có $8$ thẻ chẵn và $8$ thẻ lẻ. Không gian mẫu: $n(\Omega) = C_{16}^4 = 1820$.\\
%        Số cách chọn $4$ thẻ đều chẵn là $n(A) = C_8^4 = 70$.\\
%        Xác suất $P(A) = \frac{70}{1820} = \frac{1}{26}$.
%    }
\dotlineEX{1}
\end{ex}
\begin{ex}%[792231] [MĐ1]
    Lớp 11B có $25$ đoàn viên trong đó $10$ nam và $15$ nữ. Chọn ngẫu nhiên $3$ đoàn viên trong lớp để tham dự hội trại ngày 26 tháng 3. Tính xác suất để $3$ đoàn viên được chọn có $2$ nam và $1$ nữ.
    \choice
    {$\frac{7}{920}$}
    {\True $\frac{27}{92}$}
    {$\frac{3}{115}$}
    {$\frac{9}{92}$}
%    \loigiai{
%        Không gian mẫu: $n(\Omega) = C_{25}^3 = 2300$.\\
%        Số cách chọn $2$ nam và $1$ nữ là $n(A) = C_{10}^2 \cdot C_{15}^1 = 45 \cdot 15 = 675$.\\
%        Xác suất $P(A) = \frac{675}{2300} = \frac{27}{92}$.
%    }
\dotlineEX{1}
\end{ex}
\begin{ex}%[792246] [MĐ1]
    Một bình đựng $4$ viên bi đỏ, $3$ viên bi xanh, $2$ viên bi vàng. Lấy ngẫu nhiên từ bình ra $3$ viên bi. Tính xác suất để lấy được $3$ viên bi có đủ ba màu.
    \choice
    {$\frac{1}{3}$}
    {\True $\frac{2}{7}$}
    {$\frac{2}{3}$}
    {$\frac{1}{5}$}
%    \loigiai{
%        Tổng số bi là $9$. Không gian mẫu $n(\Omega) = C_9^3 = 84$.\\
%        Để lấy được $3$ viên bi có đủ $3$ màu, ta chọn mỗi màu $1$ viên: $n(A) = C_4^1 \cdot C_3^1 \cdot C_2^1 = 4 \cdot 3 \cdot 2 = 24$.\\
%        Xác suất $P(A) = \frac{24}{84} = \frac{2}{7}$.
%    }
\dotlineEX{1}
\end{ex}
\begin{ex}%[792247] [MĐ1]
    Có $9$ tấm thẻ được đánh số từ $1$ đến $9$. Chọn ngẫu nhiên ra $2$ tấm thẻ bất kì. Tính xác suất để tích của hai số trên $2$ tấm thẻ đã lấy là một số chẵn.
    \choice
    {$\frac{2}{3}$}
    {$\frac{1}{3}$}
    {\True $\frac{13}{18}$}
    {$\frac{5}{18}$}
    \loigiai{
        $9$ tấm thẻ gồm $5$ thẻ lẻ và $4$ thẻ chẵn. Không gian mẫu: $n(\Omega) = C_9^2 = 36$.\\
        Gọi biến cố đối $\overline{A}$: "Tích hai số là số lẻ" $\Leftrightarrow$ Cả hai thẻ lấy ra đều là số lẻ.\\
        $n(\overline{A}) = C_5^2 = 10 \Rightarrow P(\overline{A}) = \frac{10}{36} = \frac{5}{18}$.\\
        Xác suất cần tìm $P(A) = 1 - P(\overline{A}) = 1 - \frac{5}{18} = \frac{13}{18}$.
    }
\end{ex}
\begin{ex}%[792248] [MĐ2]
    Trong một lớp học gồm có $15$ học sinh nam và $10$ học sinh nữ. Giáo viên gọi ngẫu nhiên $4$ học sinh lên bảng giải bài tập. Tính xác suất để $4$ học sinh được gọi có cả nam và nữ.
    \choice
    {\True $\approx 87,5\%$}
    {$\approx 82,5\%$}
    {$\approx 80,5\%$}
    {$\approx 84,5\%$}
    \loigiai{
        Không gian mẫu: $n(\Omega) = C_{25}^4 = 12650$.\\
        Biến cố đối $\overline{A}$: "$4$ học sinh được gọi cùng giới tính".\\
        $n(\overline{A}) = C_{15}^4 + C_{10}^4 = 1365 + 210 = 1575$.\\
        Xác suất $P(\overline{A}) = \frac{1575}{12650} \approx 0,1245$.\\
        Xác suất cần tìm $P(A) = 1 - 0,1245 = 0,8755 = 87,55\% \approx 87,5\%$.
    }
\end{ex}
\begin{ex}%[792233] [MĐ2]
    Một đội ngũ giáo viên gồm $8$ thầy giáo dạy Toán học, $5$ cô giáo dạy Vật lý và $3$ cô giáo dạy Hóa học. Sở Giáo dục cần chọn ra $4$ người để chấm bài thi THPT Quốc gia, tính xác suất trong $4$ người được chọn phải có cô giáo và có đủ ba bộ môn.
    \choice
    {$\frac{5}{9}$}
    {\True $\frac{3}{7}$}
    {$\frac{4}{7}$}
    {$\frac{4}{9}$}
    \loigiai{
        Tổng số giáo viên là $8 + 5 + 3 = 16$ người.\\
        Không gian mẫu: Chọn $4$ người từ $16$ người $\Rightarrow n(\Omega) = C_{16}^4 = 1820$.\\
        Vì giáo viên môn Toán đều là thầy, môn Lý và Hóa đều là cô, nên để có đủ $3$ môn thì bắt buộc phải có giáo viên Lý và Hóa (tức là đã chắc chắn có cô giáo). Do đó ta chỉ cần đếm số cách chọn $4$ người có đủ $3$ môn. Các trường hợp gồm:
        \begin{itemize}
            \item $2$ Toán, $1$ Lý, $1$ Hóa: $C_8^2 \cdot C_5^1 \cdot C_3^1 = 28 \cdot 5 \cdot 3 = 420$ cách.
            \item $1$ Toán, $2$ Lý, $1$ Hóa: $C_8^1 \cdot C_5^2 \cdot C_3^1 = 8 \cdot 10 \cdot 3 = 240$ cách.
            \item $1$ Toán, $1$ Lý, $2$ Hóa: $C_8^1 \cdot C_5^1 \cdot C_3^2 = 8 \cdot 5 \cdot 3 = 120$ cách.
        \end{itemize}
        $n(A) = 420 + 240 + 120 = 780$.\\
        Xác suất $P(A) = \frac{780}{1820} = \frac{3}{7}$.
    }
\end{ex}
\begin{ex}%[792234] [MĐ2]
    Lấy ngẫu nhiên hai viên bi từ một thùng gồm $4$ bi xanh, $5$ bi đỏ và $6$ bi vàng. Tính xác suất để lấy được hai viên bi khác màu.
    \choice
    {$67,6\%$}
    {$29,5\%$}
    {$32,4\%$}
    {\True $70,5\%$}
    \loigiai{
        Tổng số bi trong thùng là $15$ viên. Không gian mẫu $n(\Omega) = C_{15}^2 = 105$.\\
        Gọi biến cố đối $\overline{A}$: "Lấy được $2$ viên bi cùng màu".
        $n(\overline{A}) = C_4^2 + C_5^2 + C_6^2 = 6 + 10 + 15 = 31$.\\
        Xác suất của biến cố đối: $P(\overline{A}) = \frac{31}{105}$.\\
        Xác suất cần tìm $P(A) = 1 - \frac{31}{105} = \frac{74}{105} \approx 70,476\% \approx 70,5\%$.
    }
\end{ex}
\begin{ex}%[792236] [MĐ2]
    Một hộp đựng $9$ viên bi trong đó có $4$ viên bi đỏ và $5$ viên bi xanh. Lấy ngẫu nhiên từ hộp $3$ viên bi. Tìm xác suất để $3$ viên bi lấy ra có ít nhất $2$ viên bi màu xanh.
    \choice
    {$\frac{10}{21}$}
    {$\frac{5}{14}$}
    {\True $\frac{25}{42}$}
    {$\frac{5}{42}$}
    \loigiai{
        Không gian mẫu: $n(\Omega) = C_9^3 = 84$.\\
        Có ít nhất $2$ viên bi màu xanh gồm các trường hợp:
        \begin{itemize}
            \item $2$ xanh, $1$ đỏ: $C_5^2 \cdot C_4^1 = 10 \cdot 4 = 40$ cách.
            \item $3$ xanh, $0$ đỏ: $C_5^3 \cdot C_4^0 = 10 \cdot 1 = 10$ cách.
        \end{itemize}
        Số cách thỏa mãn: $n(A) = 40 + 10 = 50$.\\
        Xác suất $P(A) = \frac{50}{84} = \frac{25}{42}$.
    }
\end{ex}
\begin{ex}%[792237] [MĐ2]
    Một tổ có $5$ học sinh nam và $6$ học sinh nữ. Giáo viên chọn ngẫu nhiên $3$ học sinh để làm trực nhật. Tính xác suất để $3$ học sinh được chọn có cả nam và nữ.
    \choice
    {$\frac{3}{8}$}
    {$\frac{24}{25}$}
    {\True $\frac{9}{11}$}
    {$\frac{3}{4}$}
    \loigiai{
        Tổng số học sinh là $11$. Không gian mẫu $n(\Omega) = C_{11}^3 = 165$.\\
        Biến cố đối $\overline{A}$: "$3$ học sinh được chọn chỉ có nam hoặc chỉ có nữ".\\
        $n(\overline{A}) = C_5^3 + C_6^3 = 10 + 20 = 30$.\\
        Xác suất $P(\overline{A}) = \frac{30}{165} = \frac{2}{11}$.\\
        Xác suất cần tìm $P(A) = 1 - P(\overline{A}) = 1 - \frac{2}{11} = \frac{9}{11}$.
    }
\end{ex}
\begin{ex}Để kiểm tra chất lượng nước uống đóng chai của ba phân xưởng X, Y, Z, đoàn kiểm tra lấy ngẫu nhiên $7$ chai nước của phân xưởng X, $4$ chai nước của phân xưởng Y và $5$ chai nước của phân xưởng Z. Sau đó lấy ngẫu nhiên $4$ chai nước từ $16$ chai ở trên để kiểm tra hàm lượng fluor theo quy định.
    \choiceTF
    {Số cách chọn $4$ chai nước từ $16$ chai của ba phân xưởng X, Y, Z để phân tích là $1960$}
    {\True Xác suất để cả ba phân xưởng X, Y, Z đều có chai nước được chọn để phân tích là $50,0\%$}
    {\True Xác suất để phân xưởng X có ít nhất một chai nước được chọn để phân tích xấp xỉ $93,1\%$}
    {Xác suất để không có chai nào của phân xưởng X được chọn là $\frac{11}{130}$}
    \loigiai{
        \begin{itemize}
            \item a) \textbf{Sai.} Tổng số chai nước là $7 + 4 + 5 = 16$. Số cách chọn $4$ chai từ $16$ chai là $C_{16}^4 = 1820$ cách (chứ không phải 1960).
            \item b) \textbf{Đúng.} Không gian mẫu: $n(\Omega) = 1820$. Số cách chọn $4$ chai có đủ $3$ phân xưởng là $C_7^2 \cdot C_4^1 \cdot C_5^1 + C_7^1 \cdot C_4^2 \cdot C_5^1 + C_7^1 \cdot C_4^1 \cdot C_5^2 = 420 + 210 + 280 = 910$. Xác suất là $P = \frac{910}{1820} = 50\%$.
            \item c) \textbf{Đúng.} Biến cố "Không có chai nào của phân xưởng X được chọn" (chỉ chọn từ Y và Z, tổng 9 chai) có số cách là $C_9^4 = 126$. Xác suất để có ít nhất 1 chai của X là $1 - \frac{126}{1820} = \frac{121}{130} \approx 93,1\%$.
            \item d) \textbf{Sai.} Xác suất để không có chai nào của phân xưởng X được chọn là $\frac{126}{1820} = \frac{9}{130}$.
        \end{itemize}
    }
\end{ex}
\begin{ex}Cho hai hộp chứa thẻ. Hộp 1 chứa các thẻ đánh số thứ tự khác nhau từ $0$ đến $100$. Hộp 2 chứa các thẻ đánh số thứ tự khác nhau từ $0$ đến $10$.
    \choiceTF
    {\True Chọn ngẫu nhiên $2$ thẻ từ hộp 1. Số cách chọn là $5050$}
    {Chọn ngẫu nhiên $2$ thẻ từ hộp 2. Xác suất để $2$ số thứ tự được chọn có tích là $1$ số lẻ là $\frac{3}{11}$}
    {\True Chọn ngẫu nhiên từ mỗi hộp $1$ thẻ. Xác suất để tổng $2$ số thứ tự không vượt quá $90$ là $\frac{86}{101}$}
    {Hộp 1 chứa $100$ thẻ}
%    \loigiai{
%        \begin{itemize}
%            \item a) \textbf{Đúng.} Hộp 1 có các số từ $0$ đến $100$ nên có $101$ thẻ. Số cách chọn $2$ thẻ từ $101$ thẻ là $C_{101}^2 = 5050$ cách.
%            \item b) \textbf{Sai.} Hộp 2 có $11$ thẻ (6 chẵn, 5 lẻ). Tích $2$ số là số lẻ khi cả hai số đều lẻ. Xác suất là $\frac{C_5^2}{C_{11}^2} = \frac{10}{55} = \frac{2}{11}$.
%            \item c) \textbf{Đúng.} Không gian mẫu là $101 \cdot 11 = 1111$. Xét biến cố đối $x+y \ge 91$, đếm được tổng cộng $10+11+ \dots +20 = 165$ cặp. Xác suất là $1 - \frac{165}{1111} = \frac{86}{101}$.
%            \item d) \textbf{Sai.} Hộp 1 chứa các thẻ đánh số từ $0$ đến $100$, do đếm cả số $0$ nên hộp 1 phải có tổng cộng $101$ thẻ.
%        \end{itemize}
%    }
\dotlineEX{1}
\end{ex}
\begin{ex}
    Cho hai đường thẳng song song $d_1, d_2$. Trên $d_1$ có $6$ điểm phân biệt được tô màu đỏ, trên $d_2$ có $4$ điểm phân biệt được tô màu xanh.
    \choiceTF
    {Số tam giác được tạo thành khi nối $10$ điểm trên lại với nhau là $120$}
    {\True Xét tất cả các đoạn thẳng được tạo thành từ $10$ điểm trên. Chọn ngẫu nhiên một đoạn thẳng, khi đó xác suất thu được đoạn thẳng có hai đỉnh khác màu là $\frac{8}{15}$}
    {\True Xét tất cả các tam giác được tạo thành khi nối các điểm đó với nhau. Chọn ngẫu nhiên một tam giác, khi đó xác suất để thu được tam giác có hai đỉnh màu đỏ là $\frac{5}{8}$}
    {Số đoạn thẳng được tạo thành từ $10$ điểm trên là $90$}
    \loigiai{
        \begin{itemize}
            \item a) \textbf{Sai.} Số tam giác tạo thành là $C_6^2 \cdot C_4^1 + C_6^1 \cdot C_4^2 = 60 + 36 = 96$ tam giác.
            \item b) \textbf{Đúng.} Số đoạn thẳng có hai đỉnh khác màu là $C_6^1 \cdot C_4^1 = 24$. Không gian mẫu là $C_{10}^2 = 45$. Xác suất là $\frac{24}{45} = \frac{8}{15}$.
            \item c) \textbf{Đúng.} Tam giác có hai đỉnh màu đỏ có số cách tạo là $C_6^2 \cdot C_4^1 = 60$. Xác suất là $\frac{60}{96} = \frac{5}{8}$.
            \item d) \textbf{Sai.} Tổng số đoạn thẳng tạo thành từ 10 điểm chỉ là $C_{10}^2 = 45$ (không phải 90).
        \end{itemize}
    }
    \dotlineEX{1}
\end{ex}
\begin{ex}
    Năm bạn Nhân, Lễ, Nghĩa, Trí và Tín xếp hàng một cách ngẫu nhiên thành một hàng ngang để chụp ảnh.
    \choiceTF
    {Số cách xếp là $60$}
    {\True Xác suất của biến cố: "Nhân và Tín không đứng cạnh nhau" bằng $\frac{3}{5}$}
    {\True Xác suất của biến cố: "Trí không đứng ở đầu hàng" bằng $\frac{4}{5}$}
    {Xác suất của biến cố: "Nhân và Tín đứng cạnh nhau" bằng $\frac{1}{5}$}
    \loigiai{
        \begin{itemize}
            \item a) \textbf{Sai.} Số cách xếp ngẫu nhiên $5$ bạn là một hoán vị của 5 phần tử: $P_5 = 5! = 120$ cách.
            \item b) \textbf{Đúng.} Số cách để Nhân và Tín đứng cạnh nhau là $2! \cdot 4! = 48$. Xác suất đứng cạnh nhau là $\frac{48}{120} = \frac{2}{5}$. Xác suất không đứng cạnh nhau là $1 - \frac{2}{5} = \frac{3}{5}$.
            \item c) \textbf{Đúng.} Trí đứng đầu hàng có $4! = 24$ cách. Xác suất là $\frac{24}{120} = \frac{1}{5}$. Xác suất Trí không đứng đầu hàng là $1 - \frac{1}{5} = \frac{4}{5}$.
            \item d) \textbf{Sai.} Xác suất để Nhân và Tín đứng cạnh nhau là $\frac{2}{5}$.
        \end{itemize}
    }
\end{ex}
\begin{ex}
    Để kiểm định chất lượng học sinh giỏi Toán trong một phường của các đội tuyển từ ba trường A, B, C. Đội trường A có 5 học sinh, đội trường B có 6 học sinh, đội trường C có 7 học sinh. Chọn ngẫu nhiên 5 bạn từ các bạn trên để kiểm tra chất lượng học tập.
    \choiceTF
    {Số cách chọn 5 bạn từ các bạn trên để kiểm tra là $8586$}
    {\True Xác suất để cả 3 trường học đều có học sinh được chọn để kiểm tra xấp xỉ $70,67\%$}
    {\True Xác suất để trường B có ít nhất 4 học sinh được chọn để kiểm tra xấp xỉ $2,17\%$}
    {Số cách chọn 5 học sinh chỉ thuộc trường C là $25$}
    \loigiai{
        \begin{itemize}
            \item a) \textbf{Sai.} Tổng số học sinh là $18$. Số cách chọn là $C_{18}^5 = 8568$ cách.
            \item b) \textbf{Đúng.} Số cách vi phạm (chỉ thuộc 1 hoặc 2 trường) là $2513$. Số cách đủ 3 trường là $8568 - 2513 = 6055$. Xác suất $P = \frac{6055}{8568} \approx 70,67\%$.
            \item c) \textbf{Đúng.} Chọn 4 học sinh B và 1 học sinh khác có $C_6^4 \cdot C_{12}^1 = 180$. Chọn 5 học sinh B có $C_6^5 = 6$. Tổng là $186$ cách. Xác suất là $\frac{186}{8568} \approx 2,17\%$.
            \item d) \textbf{Sai.} Đội trường C có 7 học sinh. Số cách chọn 5 học sinh từ 7 học sinh này là $C_7^5 = 21$ cách (không phải 25).
        \end{itemize}
    }
\end{ex}
\begin{ex}
    Có $6$ học sinh nam, $5$ học sinh nữ và $1$ thầy giáo xếp thành $1$ hàng ngang được sắp sẵn $16$ chiếc ghế để chụp kỉ yếu. Tính xác suất thầy giáo đứng giữa $1$ bạn nam và $1$ bạn nữ? (làm tròn kết quả đến hàng phần trăm)\shortans{0,25}
    \loigiai{
        Tổng số người là $6 + 5 + 1 = 12$ người xếp vào 16 ghế.\\
        Số phần tử không gian mẫu: $n(\Omega) = A_{16}^{12}$.\\
        Gọi biến cố $A$: "Thầy giáo đứng giữa 1 bạn nam và 1 bạn nữ".
        \begin{itemize}
            \item Tạo cụm Nam-Thầy-Nữ (hoặc Nữ-Thầy-Nam): có $2 \times C_6^1 \times C_5^1 = 60$ cách.
            \item Cụm này cần 3 ghế liền nhau. Chọn 3 ghế liền nhau từ 16 ghế có $16 - 3 + 1 = 14$ cách.
            \item $9$ người còn lại xếp vào $13$ ghế trống còn lại: có $A_{13}^9$ cách.
        \end{itemize}
        Số kết quả thuận lợi: $n(A) = 60 \times 14 \times A_{13}^9$.\\
        Xác suất $P(A) = \frac{60 \times 14 \times A_{13}^9}{A_{16}^{12}} = \frac{60 \times 14 \times \frac{13!}{4!}}{\frac{16!}{4!}} = \frac{60 \times 14 \times 13!}{16!} = \frac{60 \times 14}{16 \times 15 \times 14} = \frac{60}{240} = 0,25$.
        \hfill \textbf{Đáp số:} $0,25$.
    }
\end{ex}
\begin{ex}
    Có $1$ hộp chứa tất cả các thẻ là $1$ dãy số có $6$ chữ số khác nhau lập từ tập $S=\{1;2;3;4;5;6;7;8;9\}$. Rút ngẫu nhiên một thẻ từ hộp. Tính xác suất để rút được $1$ tấm thẻ có ghi các số $1,3,5,7$ sao cho hai chữ số $1$ và $3$ không đứng cạnh nhau? (làm tròn kết quả đến hàng phần trăm)\shortans{0,08}
    \loigiai{
        Số phần tử không gian mẫu: $n(\Omega) = A_9^6 = 60480$.\\
        Gọi biến cố $A$: "Thẻ rút được chứa 1,3,5,7 và 1,3 không đứng cạnh nhau".
        \begin{itemize}
            \item Chọn 2 số còn lại từ $\{2,4,6,8,9\}$: có $C_5^2 = 10$ cách.
            \item Xếp 4 số (gồm 5, 7 và 2 số vừa chọn): có $4! = 24$ cách.
            \item 4 số này tạo ra 5 vách ngăn. Xếp 1 và 3 vào 5 vách ngăn: có $A_5^2 = 20$ cách.
        \end{itemize}
        Số kết quả thuận lợi: $n(A) = 10 \times 24 \times 20 = 4800$.\\
        Xác suất $P(A) = \frac{4800}{60480} = \frac{10}{126} \approx 0,08$.
        \hfill \textbf{Đáp số:} $0,08$.
    }
\end{ex}
\begin{ex}
    Một bệnh viện dã chiến phân công ngẫu nhiên ca trực cho $3$ bác sĩ (Hòa, Thuận, Tiến) trong vòng $15$ ngày liên tiếp sao cho: Mỗi ngày chỉ có đúng $1$ bác sĩ trực, bác sĩ Hòa trực đúng $4$ ngày, bác sĩ Thuận trực đúng $4$ ngày và bác sĩ Tiến trực đúng $7$ ngày. Tính xác suất để bác sĩ Hòa không trực vào ngày đầu tiên và ngày cuối cùng của đợt $15$ ngày này? (làm tròn kết quả đến hàng phần trăm) \shortans{0,29}
    \loigiai{
        Số cách phân công ngẫu nhiên thỏa mãn số ngày trực của mỗi người (không gian mẫu):\\
        $n(\Omega) = C_{15}^4 \times C_{11}^4 \times C_7^7 = \frac{15!}{4!4!7!} = 900900$.\\
        Gọi $A$ là biến cố: "Bác sĩ Hòa không trực ngày đầu tiên và ngày cuối cùng".
        \begin{itemize}
            \item Số cách chọn 4 ngày cho BS Hòa từ 13 ngày (trừ ngày đầu và ngày cuối) là: $C_{13}^4 = 715$ cách.
            \item Còn lại 11 ngày. Số cách chọn 4 ngày cho BS Thuận là: $C_{11}^4 = 330$ cách.
            \item 7 ngày còn lại phân cho BS Tiến: $1$ cách.
        \end{itemize}
        Số kết quả thuận lợi: $n(A) = 715 \times 330 = 235950$.\\
        Xác suất $P(A) = \frac{235950}{900900} = \frac{15}{52} \approx 0,29$.
        \hfill \textbf{Đáp số:} $0,29$.
    }
    \dotlineEX{1}
\end{ex}
\begin{ex}
    Trong mặt phẳng tọa độ $Oxy$, một con robot di chuyển ngẫu nhiên từ điểm $O(0;0)$ đến điểm $B(6;5)$. Biết rằng mỗi bước đi, robot chỉ được phép đi lên trên $1$ đơn vị hoặc đi sang phải $1$ đơn vị. Tính xác suất để robot đi đến $B$ mà \textbf{không} đi qua điểm $A(2;3)$? (làm tròn kết quả đến hàng phần trăm)\shortans{0,68}
    \loigiai{
        Để đi từ $O(0;0)$ đến $B(6;5)$, robot phải đi tổng cộng $6+5=11$ bước, trong đó có 6 bước sang phải và 5 bước lên trên.\\
        Không gian mẫu (tổng số đường đi): $n(\Omega) = C_{11}^6 = 462$.\\
        Gọi biến cố $E$: "Robot không đi qua điểm $A(2;3)$". Ta tính gián tiếp qua biến cố đối $\overline{E}$: "Robot đi qua $A$".
        \begin{itemize}
            \item Số cách đi từ $O(0;0)$ đến $A(2;3)$ là $C_{2+3}^2 = C_5^2 = 10$.
            \item Số cách đi từ $A(2;3)$ đến $B(6;5)$ là $C_{4+2}^4 = C_6^4 = 15$.
        \end{itemize}
        Số đường đi qua $A$ là $n(\overline{E}) = 10 \times 15 = 150$.\\
        Số đường đi không qua $A$ là $n(E) = 462 - 150 = 312$.\\
        Xác suất $P(E) = \frac{312}{462} = \frac{52}{77} \approx 0,68$.
        \hfill \textbf{Đáp số:} $0,68$.
    }
%    \dotlineEX{1}
\end{ex}
\begin{ex}
    Có $6$ học sinh nam (trong đó có Anh và Tú) và $4$ học sinh nữ được xếp ngẫu nhiên thành một hàng ngang. Biết xác suất để không có $2$ học sinh nữ nào đứng cạnh nhau, đồng thời Anh và Tú phải đứng cạnh nhau là phân số tối giản $\frac{a}{b}$. Tính giá trị của $a+b$. \shortans{43}
    \loigiai{
        Không gian mẫu (xếp ngẫu nhiên 10 học sinh): $n(\Omega) = 10!$.\\
        Gọi biến cố $A$: "Không có 2 nữ cạnh nhau, Anh và Tú cạnh nhau".
        \begin{itemize}
            \item Gộp Anh và Tú thành cụm $X$. Hoán vị Anh và Tú: $2!$ cách.
            \item Cụm $X$ và 4 nam còn lại tạo thành 5 "khối nam". Hoán vị 5 khối này: $5!$ cách.
            \item 5 khối nam tạo ra 6 khoảng trống. Xếp 4 nữ vào 6 khoảng trống: $A_6^4$ cách.
        \end{itemize}
        Số kết quả thuận lợi: $n(A) = 2! \times 5! \times A_6^4 = 2 \times 120 \times 360 = 86400$.\\
        Xác suất $P(A) = \frac{86400}{10!} = \frac{1}{42}$.\\
        Vậy $a=1, b=42 \Rightarrow a+b = 43$.
        \hfill \textbf{Đáp số:} $43$.
    }
    \dotlineEX{1}
\end{ex}
\begin{ex}
    Một bãi đỗ xe hàng ngang có $18$ vị trí. Có $6$ chiếc xe ô tô khác nhau tiến vào bãi và đỗ ngẫu nhiên vào $6$ vị trí. Biết xác suất để giữa hai chiếc xe bất kỳ luôn có ít nhất hai vị trí trống là phân số tối giản $\frac{a}{b}$. Tính giá trị của $a+b$.\shortans{664}
    \loigiai{
        Không gian mẫu (chọn 6 vị trí có thứ tự từ 18 vị trí): $n(\Omega) = A_{18}^6 = 13366080$.\\
        Gọi biến cố $A$: "Giữa hai xe bất kỳ có ít nhất 2 vị trí trống".\\
        Cố định thứ tự 6 xe: $6!$ cách. Giữa 6 xe có 5 khoảng trống và 2 khoảng đầu mút. Gọi số vị trí trống ở các khoảng là $x_1, \dots, x_7$.\\
        Ta có: $x_1 + x_2 + \dots + x_7 = 12$ với $x_1, x_7 \ge 0$; $x_2, \dots, x_6 \ge 2$.\\
        Đặt $y_1 = x_1; y_7 = x_7; y_i = x_i - 2 \ge 0 \Rightarrow y_1 + \dots + y_7 = 2$.\\
        Số nghiệm nguyên không âm là $C_8^6 = 28$.\\
        Số kết quả thuận lợi: $n(A) = 6! \times 28 = 20160$.\\
        Xác suất $P(A) = \frac{20160}{13366080} = \frac{1}{663}$.\\
        Vậy $a=1, b=663 \Rightarrow a+b = 664$.
        \hfill \textbf{Đáp số:} $664$.
    }
\end{ex}
\begin{ex}
    Có bốn ngăn sách được đánh số thứ tự $1, 2, 3, 4$ và bảy quyển sách khác nhau. Bạn An xếp ngẫu nhiên bảy quyển sách nói trên vào bốn ngăn đó sao cho các quyển sách trong cùng một ngăn được xếp thành hàng ngang và có tính thứ tự. Biết xác suất để mỗi ngăn có ít nhất một quyển sách là phân số tối giản $\frac{a}{b}$. Tính giá trị của $a+b$.\shortans{7}
    \loigiai{
        Không gian mẫu: Xếp 7 quyển sách khác nhau vào 4 ngăn có tính thứ tự. Tương đương hoán vị 7 quyển sách ($7!$) rồi đặt 3 vách ngăn (có thể trùng hoặc ở hai đầu) vào $7+1=8$ khe. Số cách là $n(\Omega) = 7! \times C_{7+4-1}^{4-1} = 7! \times C_{10}^3 = 604800$.\\
        Gọi biến cố $A$: "Mỗi ngăn có ít nhất một quyển sách".\\
        Tương đương với việc đặt 3 vách ngăn vào 6 khoảng trống giữa các quyển sách (không trùng, không ở hai đầu): có $C_6^3$ cách.\\
        Số kết quả thuận lợi: $n(A) = 7! \times C_6^3 = 100800$.\\
        Xác suất $P(A) = \frac{100800}{604800} = \frac{1}{6}$.\\
        Vậy $a=1, b=6 \Rightarrow a+b = 7$.
        \hfill \textbf{Đáp số:} $7$.
    }
\end{ex}
\begin{ex}
    Có $4$ ngăn sách được đánh số thứ tự $1,2,3,4$ và $7$ quyển sách đôi một khác nhau. Bạn Thái xếp ngẫu nhiên hết $7$ quyển sách này vào $4$ ngăn đó sao cho ngăn nào cũng có sách. Các quyển sách trong cùng một ngăn được xếp thành hàng ngang và có tính thứ tự. Tính xác suất để ngăn thứ nhất chứa đúng $3$ quyển sách?\shortans{0,15}
    \loigiai{
        Theo kết quả bài trên, tổng số cách xếp 7 sách vào 4 ngăn có tính thứ tự sao cho ngăn nào cũng có sách (không gian mẫu mới) là $n(\Omega) = 7! \times C_6^3 = 100800$.\\
        Gọi $A$ là biến cố: "Ngăn 1 có đúng 3 quyển sách".
        \begin{itemize}
            \item Chọn 3 quyển sách cho ngăn 1 và xếp thứ tự: $A_7^3$ cách.
            \item 4 quyển sách còn lại xếp vào 3 ngăn (2, 3, 4) sao cho ngăn nào cũng có sách, có thứ tự. Xếp 4 sách thành 1 hàng ($4!$ cách), đặt 2 vách ngăn vào 3 khoảng trống giữa chúng ($C_3^2$ cách). Số cách là: $4! \times C_3^2 = 72$.
        \end{itemize}
        $n(A) = A_7^3 \times 72 = 210 \times 72 = 15120$.\\
        Xác suất $P(A) = \frac{15120}{100800} = 0,15$.
        \hfill \textbf{Đáp số:} $0,15$.
    }
    \dotlineEX{1}
\end{ex}
\begin{ex}
    Có $3$ ngăn sách được đánh số $1,2,3$ và $8$ quyển sách đôi một khác nhau. Bạn Thái xếp ngẫu nhiên hết $8$ quyển sách này vào $3$ ngăn sao cho mỗi ngăn có ít nhất một quyển sách và có tính thứ tự trong mỗi ngăn. Biết xác suất để số lượng sách ở ngăn $1$ lớn hơn ngăn $2$, và số lượng sách ở ngăn $2$ lớn hơn ngăn $3$ là phân số tối giản $\frac{a}{b}$. Tính $a+b$.\shortans{23}
    \loigiai{
        Không gian mẫu (mỗi ngăn $\ge 1$, có thứ tự): $n(\Omega) = 8! \times C_{8-1}^{3-1} = 8! \times C_7^2 = 846720$.\\
        Gọi biến cố $A$: "Số sách ngăn 1 > ngăn 2 > ngăn 3".\\
        Gọi $x_1, x_2, x_3$ là số sách ở các ngăn 1, 2, 3. Ta có $x_1 + x_2 + x_3 = 8$ và $x_1 > x_2 > x_3 \ge 1$.\\
        Các bộ $(x_1, x_2, x_3)$ thỏa mãn là: $(5, 2, 1)$ và $(4, 3, 1)$.\\
        Vì có tính thứ tự, số cách xếp cho mỗi bộ nghiệm tương đương với việc hoán vị 8 quyển sách ($8!$) rồi đặt cố định 2 vách ngăn theo tỷ lệ của bộ nghiệm.
        \begin{itemize}
            \item Trường hợp $(5, 2, 1)$: Số cách xếp là $8! = 40320$.
            \item Trường hợp $(4, 3, 1)$: Số cách xếp là $8! = 40320$.
        \end{itemize}
        Số kết quả thuận lợi $n(A) = 40320 + 40320 = 80640$.\\
        Xác suất $P(A) = \frac{80640}{846720} = \frac{2}{21}$.\\
        Vậy $a=2, b=21 \Rightarrow a+b = 23$.
        \hfill \textbf{Đáp số:} $23$.
    }
    \dotlineEX{1}
\end{ex}
\begin{ex}
    Tiến hành xếp ngẫu nhiên $10$ cuốn sách khác nhau vào $4$ ô hàng ngang sao cho tính từ trái qua phải, mỗi ngăn đều có ít nhất $1$ cuốn sách và các cuốn sách trong mỗi ngăn được xếp có thứ tự. Tính xác suất để sau khi xếp, mỗi ngăn có ít nhất $2$ cuốn sách? (làm tròn kết quả đến hàng phần trăm)\shortans{0,12}
    \loigiai{
        Không gian mẫu (chia 10 cuốn vào 4 ngăn, mỗi ngăn $\ge 1$, có thứ tự): $n(\Omega) = 10! \times C_{10-1}^{4-1} = 10! \times C_9^3 = 304819200$.\\
        Gọi biến cố $A$: "Mỗi ngăn có ít nhất 2 cuốn sách".\\
        Gọi $x_1, x_2, x_3, x_4$ là số sách trong 4 ngăn. Ta có $x_1 + x_2 + x_3 + x_4 = 10$ và $x_i \ge 2$.\\
        Đặt $y_i = x_i - 2 \ge 0 \Rightarrow y_1 + y_2 + y_3 + y_4 = 2$. Số nghiệm là $C_5^3 = 10$.\\
        Với mỗi bộ nghiệm, ta hoán vị 10 cuốn sách vào đúng vị trí tương ứng nên có $10!$ cách.\\
        Số kết quả thuận lợi $n(A) = 10 \times 10! = 36288000$.\\
        Xác suất $P(A) = \frac{36288000}{304819200} = \frac{10}{84} = \frac{5}{42} \approx 0,12$.
        \hfill \textbf{Đáp số:} $0,12$.
    }
    \dotlineEX{1}
\end{ex}
\begin{ex}
    Trong mặt phẳng tọa độ $Oxy$, một con châu chấu nhảy từ gốc tọa độ $O(0;0)$ đến điểm $B(9;4)$. Biết rằng mỗi bước, châu chấu chỉ nhảy được $1$ đơn vị theo hướng $\vec{i} = (1;0)$ hoặc $1$ đơn vị theo hướng $\vec{j} = (0;1)$ và mọi đường đi ngắn nhất từ $O$ đến $B$ đều có khả năng xảy ra như nhau. Tính xác suất để đường đi của châu chấu thỏa mãn điều kiện: cứ mỗi bước nhảy theo hướng $\vec{j}$ thì ngay sau đó phải nhảy liên tiếp ít nhất $2$ bước theo hướng $\vec{i}$. Biết xác suất là phân số tối giản $\frac{a}{b}$, tính $a+b$.\shortans{144}
    \loigiai{
        Để đi từ $O(0;0)$ đến $B(9;4)$ bằng đường ngắn nhất, châu chấu cần nhảy tổng cộng $9$ bước theo $\vec{i}$ (ký hiệu là R) và $4$ bước theo $\vec{j}$ (ký hiệu là U). Tổng số bước là $13$.\\
        Số phần tử không gian mẫu (chọn 4 vị trí cho bước U trong 13 bước): $n(\Omega) = C_{13}^4 = 715$.\\
        Gọi biến cố $A$: "Mỗi bước U luôn được tiếp nối bởi ít nhất 2 bước R".\\
        Theo điều kiện, ta "gói" 4 bước U và 8 bước R thành 4 cụm bắt buộc: $X = (U, R, R)$.\\
        Lúc này, ta đã dùng hết 4 bước U và 8 bước R, dư lại $9 - 8 = 1$ bước R.\\
        Việc xếp đường đi tương đương với việc xếp 4 cụm $X$ và 1 bước R thành một hàng ngang. Số cách xếp là số vị trí để đặt 1 bước R vào 5 vị trí trống (trước, giữa và sau 4 cụm $X$): $n(A) = C_5^1 = 5$ cách.\\
        Xác suất của biến cố là: $P(A) = \frac{5}{715} = \frac{1}{143}$.\\
        Vậy $a = 1, b = 143 \Rightarrow a+b = 144$.
        \hfill \textbf{Đáp số:} $144$.
    }
    \dotlineEX{1}
\end{ex}
\begin{ex}
    Thầy Quân có $10$ quyển vở giống nhau để phát ngẫu nhiên cho $10$ bạn học sinh khác nhau (được chia làm 3 tổ: Tổ 1 có 3 bạn, Tổ 2 có 3 bạn, Tổ 3 có 4 bạn). Mỗi bạn có thể không nhận vở hoặc nhận nhiều hơn một quyển. Tính xác suất để Tổ 1 nhận tổng số vở chẵn và chia đều cho các thành viên trong tổ, đồng thời Tổ 2 nhận tổng số vở lẻ và có đúng $2$ bạn nhận số vở bằng nhau còn bạn trực nhật nhận nhiều hơn $1$ quyển là phân số tối giản $\dfrac{a}{b}$. Tính giá trị của $b-100a$ 
    \shortans{7189}
    \loigiai{
        Số phần tử không gian mẫu (chia 10 quyển vở giống nhau cho 10 người, mỗi người nhận $x_i \ge 0$ quyển): $n(\Omega) = C_{10+10-1}^{10-1} = C_{19}^9 = 92378$.\\
        Gọi biến cố $A$: "Cách phát vở thỏa mãn các quy định của Tổ 1 và Tổ 2".
        \begin{itemize}
            \item Tổ 1 (3 bạn): Nhận số vở như nhau $k$. Tổng số vở $S_1 = 3k$. Để $S_1$ chẵn thì $k$ chẵn ($k \in \{0, 2\}$). Vậy $S_1 \in \{0, 6\}$. Có đúng $1$ cách chia nội bộ Tổ 1 cho mỗi trường hợp.
            \item Tổ 2 (3 bạn): Sự chia đều bị phá vỡ đúng 1 quyển cho bạn trực nhật, tức là có dạng $(x, x, x+1)$. Tổng số vở $S_2 = 3x+1$. Để $S_2$ lẻ thì $x$ chẵn ($x \in \{0, 2\}$).
            \begin{itemize}
                \item $x=0 \Rightarrow S_2=1$. Phân phối $(0,0,1)$ có $3$ cách chọn người nhận 1 quyển.
                \item $x=2 \Rightarrow S_2=7$. Phân phối $(2,2,3)$ có $3$ cách chọn người nhận 3 quyển.
            \end{itemize}
            \item Tổ 3 (4 bạn): Nhận $S_3 = 10 - S_1 - S_2$ quyển tùy ý. Số cách phân chia là $C_{S_3+4-1}^{4-1} = C_{S_3+3}^3$.
        \end{itemize}
        Các trường hợp xảy ra:
        \begin{itemize}
            \item TH1: $S_1 = 0, S_2 = 1 \Rightarrow S_3 = 9$. Số cách: $1 \times 3 \times C_{12}^3 = 660$.
            \item TH2: $S_1 = 0, S_2 = 7 \Rightarrow S_3 = 3$. Số cách: $1 \times 3 \times C_6^3 = 60$.
            \item TH3: $S_1 = 6, S_2 = 1 \Rightarrow S_3 = 3$. Số cách: $1 \times 3 \times C_6^3 = 60$.
        \end{itemize}
        Số kết quả thuận lợi: $n(A) = 660 + 60 + 60 = 780$.\\
        Xác suất $P(A) = \frac{780}{92378} = \frac{390}{46189}$.\\
        Vậy $a=390$, $b=46189$ và $b-100a=7189$.
    }
    \dotlineEX{2}
\end{ex}
\begin{ex}
    Trong mặt phẳng tọa độ $Oxy$, cho điểm $A(30;0)$ và điểm $B(0;30)$. Một con châu chấu nhảy ngẫu nhiên từ bên ngoài vào một điểm có tọa độ nguyên nằm bên trong hoặc nằm trên cạnh của tam giác $OAB$. Biết xác suất để con châu chấu nhảy vào điểm có tổng hoành độ và tung độ lớn hơn $18$ là phân số tối giản $\frac{a}{b}$. Tính giá trị của $a+b$.\shortans{401}
    \loigiai{
        Điểm có tọa độ nguyên $(x;y)$ thuộc tam giác $OAB$ thỏa mãn: $x \ge 0; y \ge 0; x + y \le 30$.\\
        Gọi $z \ge 0$ ($z \in \mathbb{Z}$) là "phần bù", ta có phương trình: $x + y + z = 30$.\\
        Số phần tử của không gian mẫu (số điểm thỏa mãn) chính là số nghiệm nguyên không âm của phương trình trên: $n(\Omega) = C_{30+3-1}^{3-1} = C_{32}^2 = 496$.\\
        Gọi biến cố $E$: "Điểm được chọn có tổng hoành độ và tung độ lớn hơn 18".\\
        Tức là $x + y > 18 \Leftrightarrow x + y \ge 19$. Suy ra $0 \le z \le 11$.\\
        Với mỗi giá trị $z \in \{0; 1; \dots; 11\}$ cố định, phương trình $x+y = 30-z$ có $(30-z)+1 = 31-z$ cặp nghiệm nguyên không âm $(x;y)$.\\
        Số kết quả thuận lợi là: $n(E) = \sum_{z=0}^{11} (31-z) = 31 + 30 + \dots + 20 = \frac{12 \times (31 + 20)}{2} = 306$.\\
        Xác suất cần tìm: $P(E) = \frac{306}{496} = \frac{153}{248}$.\\
        Vậy $a=153, b=248 \Rightarrow a+b = 401$.
        \hfill \textbf{Đáp số:} $401$.
    }
    \dotlineEX{1}
\end{ex}
\begin{ex}
    Tại một buổi họp, có $40$ trợ giảng đang ngồi trên một cái bàn tròn. Người quản lý chọn ngẫu nhiên $7$ bạn trợ giảng để kiểm tra chéo đáp án. Tính xác suất để trong $7$ trợ giảng được chọn có ít nhất $2$ bạn ngồi cạnh nhau? (làm tròn kết quả đến hàng phần trăm)\shortans{0,72}
    \loigiai{
        Số phần tử của không gian mẫu (chọn ngẫu nhiên 7 bạn từ 40 bạn): $n(\Omega) = C_{40}^7 = 18643560$.\\
        Gọi biến cố $A$: "Trong 7 trợ giảng được chọn có ít nhất 2 bạn ngồi cạnh nhau".\\
        Xét biến cố đối $\overline{A}$: "Không có 2 bạn nào trong 7 người được chọn ngồi cạnh nhau".\\
        Theo công thức chọn $k$ người không kề nhau trên vòng tròn $n$ người:
        $$n(\overline{A}) = \frac{n}{n-k} C_{n-k}^k = \frac{40}{40-7} C_{40-7}^7 = \frac{40}{33} C_{33}^7 = 5178240.$$
        Số kết quả thuận lợi cho $A$ là: $n(A) = n(\Omega) - n(\overline{A}) = 18643560 - 5178240 = 13465320$.\\
        Xác suất của biến cố $A$ là: $P(A) = \frac{13465320}{18643560} = \frac{112211}{155363} \approx 0,7222 \approx 0,72$.
        \hfill \textbf{Đáp số:} $0,72$.
    }
\end{ex}
\Closesolutionfile{ans}